% GENERATED FILE. Edit canonical lessons, then run npm run generate:books.
% canonical-source-hash: fnv1a64:74e46329427d1fea
% canonical-lessons: JA-C476-itoko, JA-C476-futago, JA-C476-hitorikko, JA-C476-chounan, JA-C476-choujo

\chapter{Only Children and Twins}
\label{ch:ja-only-children-and-twins}
\hlchaptermodality{476}

\begin{chapteropening}
\textbf{By the end of this chapter:} \emph{I can talk about brothers, sisters and cousins.}

Complete the last lesson of chapter 476: I can talk about brothers, sisters and cousins.
\end{chapteropening}

\section[\texorpdfstring{\ja{いとこ} (itoko)}{いとこ (itoko)}]{\ja{いとこ} (itoko) — a cousin}

\label{lesson:JA-C476-itoko}



\begin{quote}
Before the new one: say the Japanese for a small plate, then the Japanese for flavouring.
\end{quote}

\subsection*{You'll want to know: \ja{いとこ}}
\textbf{\ja{いとこ}} — \emph{itoko} — \textquotedblleft{}a cousin\textquotedblright{}.

\begin{quote}
\textbf{\ja{なつやすみに} \ja{いとこと} \ja{あそびました。}} — \emph{natsuyasumi ni itoko to asobimashita} — \textquotedblleft{}I played with my cousin in the summer holidays.\textquotedblright{}
\end{quote}

\begin{cousinweb}
\textbf{\ja{いとこ}} (\emph{itoko}) is a native Japanese word (wago), inherited rather than borrowed from Chinese; JMdict writes it in kana (JMdict).
\end{cousinweb}

\begin{grammarlens}[title={What you've built}]
One more word for describing what you lived and felt.
\end{grammarlens}

\subsection*{Your turn}
\begin{itemize}
\raggedright
  \item \emph{Say it:} \emph{itoko}
  \item \emph{Say it:} \emph{itoko}, once more
  \item \emph{Say it:} \emph{natsuyasumi ni itoko to asobimashita} — I played with my cousin in the summer holidays.
  \item \emph{Recall:} say the Japanese for a small plate, then the Japanese for flavouring, then say \emph{itoko} again
\end{itemize}

\begin{tcolorbox}[breakable,colback=teal!4,colframe=teal!35!black,title={Before you move on}]
What does \ja{いとこ} mean? (\textquotedblleft{}A cousin\textquotedblright{}.) Say it once more.
\end{tcolorbox}
\section[\texorpdfstring{\ja{ふたご} (futago)}{ふたご (futago)}]{\ja{ふたご} (futago) — twins; a twin}

\label{lesson:JA-C476-futago}



\begin{quote}
Before the new one: say the Japanese for flavouring, then the Japanese for a cousin.
\end{quote}

\subsection*{You'll want to know: \ja{ふたご}}
\textbf{\ja{ふたご}} — \emph{futago} — \textquotedblleft{}twins; a twin\textquotedblright{}.

\begin{quote}
\textbf{\ja{となりの} \ja{こどもは} \ja{ふたごです。}} — \emph{tonari no kodomo wa futago desu} — \textquotedblleft{}The children next door are twins.\textquotedblright{}
\end{quote}

\begin{cousinweb}
\textbf{\ja{ふたご}} (\emph{futago}) is a native Japanese compound (wago), inherited rather than borrowed from Chinese: \emph{futa}, \textquotedblleft{}two\textquotedblright{} and \emph{go}, \textquotedblleft{}child (ko)\textquotedblright{} (JMdict, KANJIDIC2).
\end{cousinweb}

\begin{grammarlens}[title={What you've built}]
One more word for describing what you lived and felt.
\end{grammarlens}

\subsection*{Your turn}
\begin{itemize}
\raggedright
  \item \emph{Say it:} \emph{futago}
  \item \emph{Say it:} \emph{futago}, once more
  \item \emph{Say it:} \emph{tonari no kodomo wa futago desu} — The children next door are twins.
  \item \emph{Recall:} say the Japanese for flavouring, then the Japanese for a cousin, then say \emph{futago} again
\end{itemize}

\begin{tcolorbox}[breakable,colback=teal!4,colframe=teal!35!black,title={Before you move on}]
What does \ja{ふたご} mean? (\textquotedblleft{}Twins; a twin\textquotedblright{}.) Say it once more.
\end{tcolorbox}
\section[\texorpdfstring{\ja{ひとりっこ} (hitorikko)}{ひとりっこ (hitorikko)}]{\ja{ひとりっこ} (hitorikko) — an only child}

\label{lesson:JA-C476-hitorikko}



\begin{quote}
Before the new one: say the Japanese for a cousin, then the Japanese for twins.
\end{quote}

\subsection*{You'll want to know: \ja{ひとりっこ}}
\textbf{\ja{ひとりっこ}} — \emph{hitorikko} — \textquotedblleft{}an only child\textquotedblright{}.

\begin{quote}
\textbf{\ja{かれは} \ja{ひとりっこです。}} — \emph{kare wa hitorikko desu} — \textquotedblleft{}He is an only child.\textquotedblright{}
\end{quote}

\begin{cousinweb}
\textbf{\ja{ひとりっこ}} (\emph{hitorikko}) is a native Japanese compound (wago), inherited rather than borrowed from Chinese: \emph{hitori}, \textquotedblleft{}one person\textquotedblright{} (\emph{hito}, \textquotedblleft{}one\textquotedblright{}, and \emph{ri}, the count of people), and \emph{ko}, \textquotedblleft{}child\textquotedblright{} (JMdict, KANJIDIC2).
\end{cousinweb}

\begin{grammarlens}[title={What you've built}]
One more word for describing what you lived and felt.
\end{grammarlens}

\subsection*{Your turn}
\begin{itemize}
\raggedright
  \item \emph{Say it:} \emph{hitorikko}
  \item \emph{Say it:} \emph{hitorikko}, once more
  \item \emph{Say it:} \emph{kare wa hitorikko desu} — He is an only child.
  \item \emph{Recall:} say the Japanese for a cousin, then the Japanese for twins, then say \emph{hitorikko} again
\end{itemize}

\begin{tcolorbox}[breakable,colback=teal!4,colframe=teal!35!black,title={Before you move on}]
What does \ja{ひとりっこ} mean? (\textquotedblleft{}An only child\textquotedblright{}.) Say it once more.
\end{tcolorbox}
\section[\texorpdfstring{\ja{ちょうなん} (chōnan)}{ちょうなん (chōnan)}]{\ja{ちょうなん} (chōnan) — the eldest son}

\label{lesson:JA-C476-chounan}



\begin{quote}
Before the new one: say the Japanese for twins, then the Japanese for an only child.
\end{quote}

\subsection*{You'll want to know: \ja{ちょうなん}}
\textbf{\ja{ちょうなん}} — \emph{chōnan} — \textquotedblleft{}the eldest son\textquotedblright{}.

\begin{quote}
\textbf{\ja{ちょうなんは} \ja{らいねん} \ja{けっこんします。}} — \emph{chōnan wa rainen kekkon shimasu} — \textquotedblleft{}Our eldest son is getting married next year.\textquotedblright{}
\end{quote}

\begin{cousinweb}
\textbf{\ja{ちょうなん}} (\emph{chōnan}) is Sino-Japanese (kango): two Chinese characters read with their on'yomi, readings borrowed from Chinese. They are \emph{chō}, \textquotedblleft{}senior\textquotedblright{} (Middle Chinese \emph{trjangX}) and \emph{nan}, \textquotedblleft{}male\textquotedblright{} (\emph{nom}) (KANJIDIC2; Middle Chinese in Baxter's transcription).
\end{cousinweb}

\begin{grammarlens}[title={What you've built}]
One more word for describing what you lived and felt.
\end{grammarlens}

\subsection*{Your turn}
\begin{itemize}
\raggedright
  \item \emph{Say it:} \emph{chōnan}
  \item \emph{Say it:} \emph{chōnan}, once more
  \item \emph{Say it:} \emph{chōnan wa rainen kekkon shimasu} — Our eldest son is getting married next year.
  \item \emph{Recall:} say the Japanese for twins, then the Japanese for an only child, then say \emph{chōnan} again
\end{itemize}

\begin{tcolorbox}[breakable,colback=teal!4,colframe=teal!35!black,title={Before you move on}]
What does \ja{ちょうなん} mean? (\textquotedblleft{}The eldest son\textquotedblright{}.) Say it once more.
\end{tcolorbox}
\section[\texorpdfstring{\ja{ちょうじょ} (chōjo)}{ちょうじょ (chōjo)}]{\ja{ちょうじょ} (chōjo) — the eldest daughter}

\label{lesson:JA-C476-choujo}



\begin{quote}
Before the new one: say the Japanese for an only child, then the Japanese for the eldest son.
\end{quote}

\subsection*{You'll want to know: \ja{ちょうじょ}}
\textbf{\ja{ちょうじょ}} — \emph{chōjo} — \textquotedblleft{}the eldest daughter\textquotedblright{}.

\begin{quote}
\textbf{\ja{ちょうじょは} \ja{いしゃに} \ja{なりました。}} — \emph{chōjo wa isha ni narimashita} — \textquotedblleft{}Our eldest daughter became a doctor.\textquotedblright{}
\end{quote}

\begin{cousinweb}
\textbf{\ja{ちょうじょ}} (\emph{chōjo}) is Sino-Japanese (kango): two Chinese characters read with their on'yomi, readings borrowed from Chinese. They are \emph{chō}, \textquotedblleft{}senior\textquotedblright{} (Middle Chinese \emph{trjangX}) and \emph{jo}, \textquotedblleft{}woman\textquotedblright{} (\emph{nrjoX}) (KANJIDIC2; Middle Chinese in Baxter's transcription).
\end{cousinweb}

\begin{grammarlens}[title={What you've built}]
One more word for describing what you lived and felt.
\end{grammarlens}

\subsection*{Your turn}
\begin{itemize}
\raggedright
  \item \emph{Say it:} \emph{chōjo}
  \item \emph{Say it:} \emph{chōjo}, once more
  \item \emph{Say it:} \emph{chōjo wa isha ni narimashita} — Our eldest daughter became a doctor.
  \item \emph{Recall:} say the Japanese for an only child, then the Japanese for the eldest son, then say \emph{chōjo} again
\end{itemize}

\begin{tcolorbox}[breakable,colback=teal!4,colframe=teal!35!black,title={Before you move on}]
What does \ja{ちょうじょ} mean? (\textquotedblleft{}The eldest daughter\textquotedblright{}.) Say it once more.
\end{tcolorbox}
